\subsection{Результаты: влияние признаков на прогноз}

В данном разделе результаты рассматриваются не как соревнование моделей, а как проверка влияния разных групп признаков на прогноз направления цены. Основной вопрос состоит в том, дают ли FinBERT-generated features самостоятельный сигнал и улучшают ли они качество прогноза при объединении с рыночными признаками.

Основным экспериментом является RUN 03. В нём используется хронологическое разбиение \texttt{chronological\_row\_balanced\_by\_decision\_date}: train, validation и test не пересекаются по датам, а validation и test имеют достаточный размер для первичного сравнения. В тестовой выборке 222 наблюдения.

\begin{table}[h]
\centering
\begin{tabular}{llrrrr}
\toprule
Model & Feature group & Accuracy & Balanced acc. & F1 & ROC-AUC \\
\midrule
Majority & none & 0.532 & 0.500 & 0.694 & 0.500 \\
Logistic Regression & prices-only & 0.410 & 0.400 & 0.502 & 0.379 \\
Logistic Regression & FinBERT-only & 0.550 & 0.529 & 0.667 & 0.512 \\
Logistic Regression & prices + FinBERT & 0.486 & 0.474 & 0.584 & 0.427 \\
LSTM & prices + FinBERT & 0.432 & 0.445 & 0.315 & 0.449 \\
\bottomrule
\end{tabular}
\caption{Качество моделей на test subset в RUN 03}
\end{table}

Первое ключевое сравнение --- \texttt{FinBERT-only} против majority baseline. Оно показывает, содержат ли текстовые embedding-признаки самостоятельную информацию о направлении цены. FinBERT-only модель показывает accuracy 0.550, balanced accuracy 0.529, F1 0.667 и ROC-AUC 0.512. По accuracy и balanced accuracy она немного превосходит majority baseline, у которого accuracy равна 0.532. Это означает, что текстовые признаки могут содержать некоторый самостоятельный сигнал. Однако величина преимущества мала, а ROC-AUC остаётся почти равным 0.5. Следовательно, в текущем эксперименте FinBERT-признаки не являются полностью бесполезными, но их самостоятельная предсказательная сила выражена слабо.

Второе и главное сравнение --- \texttt{prices + FinBERT} против \texttt{prices-only}. Оно проверяет добавочную предсказательную силу текстовых признаков поверх рыночной информации. Модель на рыночных признаках показала accuracy 0.410, balanced accuracy 0.400, F1 0.502 и ROC-AUC 0.379. Модель на объединённых признаках показала accuracy 0.486, balanced accuracy 0.474, F1 0.584 и ROC-AUC 0.427. По сравнению с \texttt{prices-only} объединённая модель показывает более высокие значения основных метрик. Это означает, что добавление FinBERT-признаков улучшает результат относительно модели, использующей только рыночные признаки.

Однако это улучшение не является достаточным для сильного вывода о добавочной силе текстов. Во-первых, \texttt{prices + FinBERT} всё ещё уступает majority baseline по accuracy. Во-вторых, ROC-AUC объединённой модели остаётся ниже 0.5. В-третьих, объединённая модель работает хуже, чем \texttt{FinBERT-only}, что означает отсутствие устойчивого выигрыша от комбинации рыночных и текстовых признаков.

Если рыночные и текстовые признаки несут комплементарную информацию, то модель \texttt{prices + FinBERT} должна работать лучше обеих отдельных групп. В текущем эксперименте этого не происходит. \texttt{FinBERT-only} показывает accuracy 0.550 и ROC-AUC 0.512, тогда как \texttt{prices + FinBERT} показывает accuracy 0.486 и ROC-AUC 0.427. Это значит, что добавление рыночных признаков к FinBERT embeddings не улучшает качество прогноза. С точки зрения исследования признаков это означает, что комплементарность рыночной и текстовой информации в текущем протоколе не подтверждается.

LSTM-модель использовалась для проверки того, помогает ли последовательная архитектура извлекать сигнал из объединённых признаков \texttt{prices + FinBERT}. Она показала accuracy 0.432, balanced accuracy 0.445, F1 0.315 и ROC-AUC 0.449. Эти значения ниже majority baseline и ниже логистической регрессии на \texttt{FinBERT-only}. Это означает, что усложнение модели не улучшило использование признаков в рамках текущего MVP-протокола.

Итоговая интерпретация состоит в следующем. FinBERT-only признаки дают небольшой прирост accuracy относительно majority baseline, что указывает на наличие слабого самостоятельного текстового сигнала. Добавление FinBERT-признаков к рыночным признакам улучшает результат относительно \texttt{prices-only}, но объединённая модель всё равно не превосходит majority baseline и уступает FinBERT-only. Комплементарность рыночных и текстовых признаков в текущем протоколе не подтверждается. Следовательно, в рамках выбранного актива, периода, горизонта прогноза и экспериментального протокола нельзя утверждать, что LLM-generated features обладают устойчивой добавочной предсказательной силой.
